\begin{frame}{Предполагаемая научная новизна}
\begin{enumerate}
\item Способ сопряжения временного модуля, предобученного на коннектомно-ограниченной нейродинамике, с биофизическим оператором ECAP: переносимая часть параметров, способы её адаптации и условия, при которых начальная параметризация влияет на точность.
\item Разграничение вклада структуры связей, обученной динамики и общего предобучения на сопоставимых контрольных моделях и простой модели кривой роста.
\item Условия применимости и пределы неопределённости моделирования кривой роста ECAP.
\end{enumerate}
\SlideGap
Сочетание известных компонентов само по себе не считается новизной. Поиск ближайших решений выявил только частичные аналоги; приоритет не заявляется, а оценка вклада зависит от результатов сравнений.
\end{frame}
